% ═══════════════════════════════════════════════════════════════
%  chapters/ch06-trigger.tex
%  Триггер: как поймать начало сигнала
% ═══════════════════════════════════════════════════════════════

\chapter{Триггер: как поймать начало сигнала}
\label{ch:trigger}

\epigraph{<<Без триггера осциллограф~--- это просто телевизор
со сломанной синхронизацией. Картинка есть, но она
бесконечно бежит по экрану>>.}{}

\section{Введение}
\label{sec:trigger-intro}

В \texttt{ch05} мы научились \emph{захватывать} сигнал с
помощью ADC и DMA. Мы получаем \emph{числа}~--- тысячи
чисел в секунду. Но \emph{как} их \emph{показать}
на экране?

\textbf{Проблема}: если просто \emph{рисовать} всё, что
пришло~--- картинка \emph{поплывёт}. Каждый \emph{новый}
кадр будет начинаться \emph{в} \emph{случайном} месте.
Синус будет \emph{дёргаться} влево-вправо.

\begin{ideabox}
\textbf{Аналогия: старый телевизор.} Помните, когда
антенна \emph{плохо} ловила~--- картинка \emph{бежала}
по экрану \emph{сверху} вниз или \emph{слева} направо.
Это~--- \emph{отсутствие} синхронизации. Триггер~---
это \emph{синхронизация} для осциллографа.
\end{ideabox}

\textbf{Триггер} (trigger)~--- это \emph{условие}, при
котором осциллограф \emph{начинает} запись. Он говорит:
<<\emph{начинай} захват, \emph{когда} сигнал \emph{пересечёт}
\emph{вот} \emph{этот} уровень \emph{вот} \emph{в} этом
направлении>>.

\subsection{Что вы узнаете}

\begin{itemize}
    \item Что такое \emph{уровень} (level) триггера.
    \item Что такое \emph{фронт} (edge)~--- по возрастанию
          или по убыванию.
    \item Какие бывают \emph{режимы} триггера.
    \item Что такое \emph{предтриггер} (pre-trigger) и зачем он нужен.
    \item Как \emph{реализовать} триггер на STM32.
\end{itemize}

\section{Что такое триггер}
\label{sec:trigger-what}

\textbf{Триггер}~--- это \emph{событие}, которое
\emph{запускает} запись. \emph{Пока} событие \emph{не}
произошло~--- осциллограф \emph{ждёт}. \emph{Когда}
произошло~--- \emph{захватывает} кадр и \emph{показывает}.

\subsection{Уровень триггера}

\textbf{Уровень} (level)~--- это \emph{напряжение}, при
котором \emph{срабатывает} триггер.

\textbf{Пример}: уровень \SI{1.5}{\volt}. Осциллограф
\emph{ждёт}, пока сигнал \emph{не} \emph{пересечёт}
\SI{1.5}{\volt}. Как только \emph{пересечёт}~--- \emph{начинает}
запись.

\begin{keybox}
\textbf{Зачем нужен уровень?} Чтобы \emph{всегда}
начинать \emph{с} \emph{одного} и \emph{того} же
места сигнала. Если сигнал~--- \emph{синус}, и мы
\emph{ждём} \SI{0}{\volt}~--- кадр \emph{начнётся}
с \emph{нуля}. \emph{Каждый} раз. \emph{Картинка}~---
\emph{стабильна}.
\end{keybox}

\subsection{Фронт триггера}

\textbf{Фронт} (edge)~--- это \emph{направление}
пересечения уровня:
\begin{itemize}
    \item \textbf{По возрастанию} (rising)~--- сигнал
          \emph{растёт} и \emph{пересекает} уровень \emph{снизу} вверх.
    \item \textbf{По убыванию} (falling)~--- сигнал
          \emph{падает} и \emph{пересекает} уровень \emph{сверху} вниз.
\end{itemize}

\begin{figure}[htbp]
\centering
\adjustbox{max width=\textwidth}{%
\begin{tikzpicture}[
    x=1cm, y=1cm,
    axis/.style={-{Latex[length=3mm]}, thick},
    signal/.style={very thick, color=clrLogicData},
    level/.style={thick, dashed, color=clrWarn},
    marker/.style={circle, fill=clrWarn, inner sep=2pt},
    label/.style={font=\footnotesize, inner sep=2pt}]
% ─── Координатная сетка ────────────────────────────
\draw[step=1cm, gray!20, very thin] (-0.5, -2) grid (10, 4);
% Оси
\draw[axis] (-0.5, 0) -- (9.5, 0) node[right] {\(t\)};
\draw[axis] (0, -1.5) -- (0, 3.5) node[above] {\(U\)};
% Уровень триггера
\draw[level] (-0.5, 1.5) -- (9.5, 1.5);
\node[label, right] at (9.5, 1.5) {уровень};
% Сигнал (синус с шумом и смещением)
\draw[signal, domain=0:9, samples=200]
    plot (\x, {1.5 + 1.2 * sin(deg(1.5 * \x))});
% Точки срабатывания триггера (по возрастанию)
\node[marker] at (1.05, 1.5) {};
\node[marker] at (5.23, 1.5) {};
\node[marker] at (9.41, 1.5) {};
% Подписи
\node[label, above] at (1.05, 1.7) {rising};
\node[label, above] at (5.23, 1.7) {rising};
\node[label, above] at (9.41, 1.7) {rising};
% Метки осей (для отладки)
\foreach \x in {0,...,9} {
    \node[font=\tiny, gray] at (\x, -2.3) {\x};}
\foreach \y in {-2,-1,0,1,2,3} {
    \node[font=\tiny, gray] at (-0.7, \y) {\y};}
\end{tikzpicture}
}% end adjustbox
\caption{Уровень триггера и точки срабатывания.
Красная пунктирная линия~--- уровень. Кружки~--- моменты,
когда сигнал \emph{пересекает} уровень \emph{по возрастанию}.}
\label{fig:trigger-level}
\end{figure}

\begin{warnbox}
\textbf{Почему именно \emph{точка} пересечения, а не \emph{весь}
сигнал?} Осциллограф \emph{не} может \emph{захватить}
\emph{всё}~--- слишком \emph{много} данных. Он \emph{захватывает}
\emph{окно} (например, \num{1024} семпла) \emph{вокруг}
точки \emph{пересечения}. Триггер \emph{говорит}, \emph{где}
начать окно.
\end{warnbox}

\section{Режимы триггера}
\label{sec:trigger-modes}

\textbf{Триггер} может работать в \emph{трёх} режимах:

\begin{itemize}
    \item \texttt{off}~--- \emph{выключен}. Осциллограф
          \emph{захватывает} всё подряд. \emph{Картинка}
          \emph{плывёт}.
    \item \texttt{normal}~--- \emph{ждёт} \emph{события}.
          Если \emph{событие} не произошло~--- \emph{ничего}
          не показывает. \emph{Картинка} \emph{стабильна},
          но \emph{если} сигнал \emph{пропал}~--- \emph{пусто}.
    \item \texttt{auto}~--- \emph{ждёт} \emph{события},
          но \emph{если} \emph{долго} \emph{нет}~---
          \emph{показывает} \emph{что} есть. \emph{Компромисс}.
\end{itemize}

\begin{keybox}
\textbf{Что выбрать?} Для \emph{стабильного} сигнала~---
\texttt{normal}. Для \emph{редких} импульсов~---
\texttt{normal} (ждём). Для \emph{поиска} сигнала~---
\texttt{auto}. \emph{По умолчанию}~--- \texttt{auto}.
\end{keybox}

\begin{table}[htbp]
\centering
\caption{Режимы триггера}
\label{tab:trigger-modes}
\small
\begin{tabularx}{\textwidth}{@{}l L L@{}}
\toprule
\textbf{Режим} & \textbf{Что делает} & \textbf{Когда использовать} \\
\midrule
\texttt{off}    & Захватывает всё подряд & Отладка, поиск \\
\texttt{normal} & Ждёт события, иначе~--- пусто & Стабильный сигнал \\
\texttt{auto}   & Ждёт, но через таймаут~--- показывает & По умолчанию \\
\bottomrule
\end{tabularx}
\end{table}

\section{Предтриггер}
\label{sec:trigger-pretrigger}

\textbf{Проблема}: если триггер сработал~--- мы \emph{начинаем}
запись. Но \emph{что} было \emph{до} триггера? \emph{Самое}
интересное~--- \emph{перед} событием (например, \emph{фронт}
импульса).

\textbf{Решение}: \textbf{предтриггер} (pre-trigger).
Буфер \emph{делится} на \emph{две} части:
\begin{itemize}
    \item \textbf{До триггера} (\texttt{pre})~--- \emph{уже}
          \emph{захвачено} \emph{до} события.
    \item \textbf{После триггера} (\texttt{post})~--- \emph{захватывается}
          \emph{после} события.
\end{itemize}

\begin{figure}[htbp]
\centering
\adjustbox{max width=\textwidth}{%
\begin{tikzpicture}[
    x=1cm, y=1cm,
    buf/.style={
        draw, thick,
        minimum height=1.5cm,
        align=center,
        font=\footnotesize
    },
    pre/.style={buf, fill=clrLogicData!20},
    post/.style={buf, fill=clrWarn!20},
    trigger/.style={
        draw, thick, fill=clrWarn,
        minimum width=0.1cm,
        minimum height=1.5cm
    },
    label/.style={font=\footnotesize, inner sep=2pt}
]
% Координатная сетка
\draw[step=1cm, gray!20, very thin] (-1, -2) grid (9, 2);
% Буфер
\node[pre,  minimum width=4cm] (pre)  at (1.5, 0) {pre (до)};
\node[post, minimum width=4cm] (post) at (6.5, 0) {post (после)};
% Триггер
\node[trigger] at (4, 0) {};
% Стрелка времени
\draw[-{Latex[length=3mm]}, thick] (-1.0, -1.5) -- (9, -1.5) node[right] {\(t\)};
% Подписи
\node[label, above] at (1.5, 0.8) {\texttt{pre\_pct}};
\node[label, below] at (1.5, -0.8) {\texttt{pre = 50\%}};
\node[label, above] at (6.5, 0.8) {\texttt{post}};
\node[label, below] at (6.5, -0.8) {\texttt{post = 50\%}};
\node[label, above] at (4, 0.8) {\textbf{триггер}};
% Метки осей
\foreach \x in {-1,...,9} {
    \node[font=\tiny, gray] at (\x, -2.3) {\x};}
\foreach \y in {-2,-1,0,1,2} {
    \node[font=\tiny, gray] at (-1.3, \y) {\y};}
\end{tikzpicture}
}% end adjustbox
\caption{Предтриггер. Буфер делится на две части:
\emph{до} триггера и \emph{после}. Триггер~--- в центре.}
\label{fig:trigger-pretrigger}
\end{figure}

\textbf{Как это работает}:
\begin{enumerate}
    \item DMA \emph{пишет} в \emph{буфер}. \emph{До} триггера~---
          \emph{все} данные \emph{идут} в \emph{pre}.
    \item Триггер \emph{сработал}. \emph{С} этого момента~---
          данные \emph{идут} в \emph{post}.
    \item Когда \emph{post} \emph{заполнен}~--- \emph{кадр}
          \emph{готов}. \emph{Отправляем} на ПК.
\end{enumerate}

\begin{keybox}
\textbf{Зачем \texttt{pre\_pct}?} Это \emph{процент} буфера,
который \emph{до} триггера. \texttt{50\%}~--- триггер
\emph{в} \emph{центре}. \texttt{25\%}~--- \emph{больше}
\emph{после}. \texttt{75\%}~--- \emph{больше} \emph{до}.
\emph{Пользователь} \emph{выбирает}.
\end{keybox}

\section{Реализация на STM32}
\label{sec:trigger-impl}

\textbf{Как} \emph{реализовать} триггер на STM32? \emph{Два}
варианта:

\subsection{Вариант 1: программный триггер}

\textbf{В} \emph{обработчике} DMA~\--- \emph{проверяем}
каждый семпл:
\begin{itemize}
    \item Если \emph{сигнал} \emph{пересёк} уровень \emph{в}
          \emph{нужном} направлении~--- \emph{запоминаем}
          позицию.
\end{itemize}

\textbf{Плюс}: \emph{просто}.
\textbf{Минус}: \emph{медленно} (\emph{много} проверок).

\subsection{Вариант 2: аппаратный триггер (Analog Watchdog)}

\textbf{ADC} \emph{умеет} \emph{сам} \emph{следить} за
уровнем. Это~--- \textbf{Analog Watchdog} (AWD).
\emph{Настраиваем} \emph{порог}~--- ADC \emph{генерирует}
\emph{прерывание}, когда \emph{сигнал} \emph{пересекает}
его.

\textbf{Плюс}: \emph{быстро}, \emph{CPU} \emph{не} \emph{тратит}
время.
\textbf{Минус}: \emph{сложнее} \emph{настроить}.

\begin{ideabox}
\textbf{Какой вариант выбрать?} Для \emph{учебного}
осциллографа~--- \emph{оба} \emph{подходят}. \emph{Мы}
начнём с \emph{программного}~--- он \emph{понятнее}.
\emph{Позже}~--- \emph{аппаратный} (\emph{если} \emph{нужна}
\emph{скорость}).
\end{ideabox}


%%%%% ======================================



\section{Код: trigger.c}
\label{sec:trigger-code}

\textbf{Теперь}~--- \emph{код}. Он \emph{короткий}, но
\emph{плотный}. \emph{Каждая} строка \emph{важна}.

\subsection{Программный триггер}

\begin{listing}[htbp]
\begin{minted}[frame=none, bgcolor=gray!5]{c}
static uint16_t trigger_level = 2048;
static int8_t   trigger_edge  = 1;  // 1 = rising, -1 = falling
static bool     trigger_found = false;
static uint16_t trigger_pos   = 0;

void trigger_check(const uint16_t *buf, uint16_t len)
{
    if (trigger_found) return;

    for (uint16_t i = 1; i < len; i++) {
        // Проверяем пересечение уровня
        bool rising = (buf[i - 1] < trigger_level &&
                       buf[i]     >= trigger_level);
        bool falling = (buf[i - 1] > trigger_level &&
                        buf[i]     <= trigger_level);

        if ((trigger_edge == 1 && rising) ||
            (trigger_edge == -1 && falling)) {
            trigger_found = true;
            trigger_pos   = i;
            break;
        }
    }
}
\end{minted}
\caption{\texttt{trigger\_check()}~--- программный триггер}
\label{lst:trigger-check}
\end{listing}

\textbf{Что здесь важно}:
\begin{itemize}
    \item \texttt{trigger\_level}~--- \emph{уровень}
          (\num{2048} = \SI{1.65}{\volt}).
    \item \texttt{trigger\_edge}~--- \emph{направление}
          (\texttt{1}~--- \emph{вверх}, \texttt{-1}~--- \emph{вниз}).
    \item \texttt{trigger\_found}~--- \emph{флаг} (нашли
          или \emph{нет}).
    \item \texttt{trigger\_pos}~--- \emph{где} нашли
          (индекс в буфере).
\end{itemize}

\begin{warnbox}
\textbf{Почему \texttt{i = 1}, а не \texttt{0}?} Потому
что \emph{сравниваем} \emph{текущий} с \emph{предыдущим}.
\texttt{buf[0]}~--- \emph{не} с \emph{чем} сравнивать.
\end{warnbox}

\subsection{Интеграция с DMA}

\begin{listing}[htbp]
\begin{minted}[frame=none, bgcolor=gray!5]{c}
void HAL_ADC_ConvHalfCpltCallback(ADC_HandleTypeDef *hadc)
{
    trigger_check(&dma_buf[0], OSC_DMA_POINTS);
    if (trigger_found) {
        push_block(&dma_buf[0], OSC_DMA_POINTS);
    }
}

void HAL_ADC_ConvCpltCallback(ADC_HandleTypeDef *hadc)
{
    trigger_check(&dma_buf[OSC_DMA_POINTS], OSC_DMA_POINTS);
    if (trigger_found) {
        push_block(&dma_buf[OSC_DMA_POINTS], OSC_DMA_POINTS);
    }
}
\end{minted}
\caption{Интеграция триггера с DMA}
\label{lst:trigger-dma}
\end{listing}

\textbf{Как это работает}:
\begin{enumerate}
    \item DMA \emph{заполняет} половину буфера.
    \item \emph{Срабатывает} прерывание.
    \item \texttt{trigger\_check()} \emph{ищет} уровень.
    \item Если \emph{нашли}~--- \texttt{push\_block()} \emph{отправляет}
          данные на ПК.
\end{enumerate}

\begin{keybox}
\textbf{Почему в прерывании?} Потому что \emph{данные}
\emph{готовы} \emph{сразу}. Если \emph{ждать} до
\emph{основного} цикла~--- \emph{DMA} \emph{перезапишет}
буфер. \emph{Прерывание}~--- \emph{единственный} способ
\emph{успеть}.
\end{keybox}

\subsection{Предтриггер}

\begin{listing}[htbp]
\begin{minted}[frame=none, bgcolor=gray!5]{c}
// Внутри push_block
static uint16_t pre_cnt = 0;
static uint16_t post_cnt = 0;
static uint16_t pre_limit = 0;

void push_block(const uint16_t *src, uint16_t count)
{
    if (!trigger_found) {
        // До триггера: копим в pre
        pre_cnt = (pre_cnt + count > pre_limit) ?
                  pre_limit : pre_cnt + count;
        // ... копирование в pre-буфер ...
    } else {
        // После триггера: копим в post
        post_cnt += count;
        // ... копирование в post-буфер ...
        if (post_cnt >= OSC_FRAME_POINTS - pre_limit) {
            // Кадр готов!
            send_frame();
            pre_cnt = post_cnt = 0;
            trigger_found = false;
        }
    }
}
\end{minted}
\caption{Предтриггер в \texttt{push\_block()}}
\label{lst:trigger-pretrigger}
\end{listing}

\textbf{Что здесь важно}:
\begin{itemize}
    \item \texttt{pre\_cnt}~--- \emph{сколько} уже \emph{накоплено}
          \emph{до} триггера.
    \item \texttt{pre\_limit}~--- \emph{максимум} \emph{до}
          (\texttt{pre\_pct} \emph{от} \emph{размера} \emph{кадра}).
    \item \texttt{post\_cnt}~--- \emph{сколько} \emph{после}.
    \item \texttt{trigger\_found}~--- \emph{флаг}.
\end{itemize}

\begin{keybox}
\textbf{Ключевая идея}: \emph{до} триггера~--- \emph{копим}
в \emph{pre}. \emph{После}~--- \emph{копим} в \emph{post}.
\emph{Когда} \emph{post} \emph{заполнен}~--- \emph{кадр}
\emph{готов}. \emph{Отправляем}.
\end{keybox}

\section{Уроки главы}
\label{sec:trigger-lessons}

\begin{enumerate}
    \item \textbf{Триггер~--- это синхронизация.} Без
          него~--- \emph{картинка} \emph{плывёт}. Он
      \emph{начинает} захват \emph{в} \emph{одном} и \emph{том}
      же месте сигнала.
    \item \textbf{Уровень}~--- \emph{напряжение}, при котором
          \emph{срабатывает} триггер.
    \item \textbf{Фронт}~--- \emph{направление} пересечения
          (\emph{вверх} или \emph{вниз}).
    \item \textbf{Режимы}: \texttt{off}, \texttt{normal},
          \texttt{auto}.  \emph{По умолчанию}~--- \texttt{auto}.
    \item \textbf{Предтриггер}~--- \emph{буфер} \emph{до}
          и \emph{после} события.
    \item \textbf{Программный} \emph{триггер}~--- \emph{просто},
          но \emph{медленно}.
    \item \textbf{Аппаратный триггер} (\emph{AWD})~--- \emph{быстрее},
          но \emph{сложнее}. 
    \item \textbf{Интеграция с DMA}~--- \emph{только} в
          \emph{прерывании}.
\end{enumerate}


\section{Чек-лист грабель}
\label{sec:trigger-pitfalls}

\begin{table}[htbp]
\centering
\caption{Грабли триггера}
\label{tab:trigger-pitfalls}
\small
\begin{tabularx}{\textwidth}{@{}c L L@{}}
\toprule
\textbf{№} & \textbf{Грабля} & \textbf{Решение} \\
\midrule
1 & Триггер не срабатывает & Проверить уровень \\
2 & Триггер срабатывает \emph{слишком} часто & Проверить фронт \\
3 & Картинка \emph{плывёт} & Включить триггер \\
4 & Триггер \emph{теряет} данные & Обработка в \emph{прерывании} \\
5 & Нет предтриггера & Делить буфер на \emph{pre} и \emph{post} \\
6 & \texttt{pre\_pct} > \num{100} & Ограничить \\
7 & Долгий \texttt{trigger\_check()} & Оптимизировать или \emph{AWD} \\
8 & Нет \emph{сброса} флага & \texttt{trigger\_found = false} \\
\bottomrule
\end{tabularx}
\end{table}

\section{Упражнения}
\label{sec:trigger-exercises}

\begin{exercisebox}
\textbf{Упражнение 1.} Нарисуйте \emph{сигнал} и \emph{уровень}
триггера. Покажите \emph{точки} срабатывания для
\emph{rising} и \emph{falling}.
\end{exercisebox}

\begin{exercisebox}
\textbf{Упражнение 2.} Что произойдёт, если \texttt{trigger\_edge = 0}?
Приведите пример, когда это \emph{полезно}.
\end{exercisebox}

\needspace{5\baselineskip}
\begin{exercisebox}
\textbf{Упражнение 3.} Реализуйте \emph{аппаратный} триггер
через \emph{AWD}. Какие регистры \emph{настроить}?
\end{exercisebox}

\begin{exercisebox}
\textbf{Упражнение 4.} Предложите способ \emph{автоматически}
выбирать \emph{уровень} триггера (например, \emph{середина}
сигнала). Как это \emph{упростит} жизнь пользователю?
\end{exercisebox}

\section{Проверьте сами}
\label{sec:trigger-check}

\begin{enumerate}
    \item Подключите Nucleo к компьютеру.
    \item Соберите прошивку (см. главу 2).
    \item Подключите генератор к \texttt{PA0} (\SI{1}{\kilo\hertz},
          \SI{1}{\volt}pp).
    \item Запустите OpenOCD + GDB.
    \item Поставьте точку останова в \texttt{trigger\_check()}.
    \item Посмотрите, \emph{как} сигнал \emph{пересекает} уровень.
    \item Сделайте \emph{скриншот} — \emph{пришлите} в книгу.
\end{enumerate}

\begin{ideabox}
\textbf{Если получилось~--- вы поняли триггер.} Это
\emph{главное}. Без триггера~--- \emph{нет} осциллографа.
\emph{Есть} только \emph{мусор} на экране. Дальше~---
\emph{кольцо кадров}, \emph{USB}, \emph{графика}.
\emph{Всё} \emph{впереди}.
\end{ideabox}


